\section{Prompt Templates and Inference Configuration}
\label{app:prompts}

This appendix documents the schemas, prompt templates, and inference settings used to generate LLM outputs. Placeholders \verb+{{TEXT}}+, \verb+{{TRANSCRIPT}}+, \verb+{{THUMBNAIL}}+, \verb+{{VIDEO\FRAMES}}+, and \verb+{{CONCEPT_DEFINITION}}+ were filled dynamically for each query.

\subsection{Structured Output Schema}
\label{app:schema}

Models were instructed to emit a \texttt{JSON} object conforming to a predefined schema with a single integer outcome in: \texttt{0}~=~``no'', \texttt{1}~=~``yes'', \texttt{2}~=~``inconclusive\_evidence'', \texttt{3}~=~``inconclusive\_definition'', which are validated against a fixed schema; invalid (malformed) outputs trigger up to three deterministic retries.\\
The \texttt{ConceptClassification} object contains the final \texttt{outcome}, optional \verb+reasoning_steps+, self-reported confidence \verb+p_correct+, and a confidence \texttt{band}. \\
For reproducibility, the overall \verb+LLMClassification+ object and \verb+ReasoningStep+ are shown below.

\begingroup\small
\begin{verbatim}
class LLMClassification(BaseModel):
    classifications: Dict[ConceptEnum, ConceptClassification]

class ConceptClassification(BaseModel):
    outcome: Optional[OutcomeEnum]
    reasoning_steps: Optional[List[ReasoningStep]]
    p_correct: Optional[int]
    band: Optional[BandEnum]

class ReasoningStep(BaseModel):
    step_number: int
    description: str
\end{verbatim}
\endgroup

\myparagraph{Schema Clarifications}
All CoT prompts enforced exactly three \verb+reasoning_steps+. We used regular expressions to normalize the outcome tokens\\
\texttt{(yes|no|inconclusive\_(definition|evidence))} for consistent parsing. The \verb+p_correct+ field was snapped to the nearest multiple of 5 in \([0,100]\), and the \texttt{band} was validated against \verb+{VL, L, M, H, VH}+.

\subsection{Representative Prompt Templates}
\label{app:prompts-templates}

\myparagraph{Text-only Baseline (System Prompt)}
\begin{verbatim}
# ROLE AND GOAL
You are a meticulous Chief Marketing Officer (CMO) ...
...
\end{verbatim}

\myparagraph{Text-only User Prompt}
\begin{verbatim}
Please classify the following content:
--- CONTENT START ---
{{TEXT}}
--- CONTENT END ---
\end{verbatim}

\myparagraph{Multimodal User Prompt}
\begin{verbatim}
Analyze the following multimodal content.

--- START OF MULTIMODAL CONTENT ---
<VIDEO\FRAMES> {{VIDEO\FRAMES}} </VIDEO\FRAMES>
<THUMBNAIL> {{THUMBNAIL}} </THUMBNAIL>
<TRANSCRIPT> {{TRANSCRIPT}} </TRANSCRIPT>
<CONTENT_TEXT> {{TEXT}} </CONTENT_TEXT>
--- END OF MULTIMODAL CONTENT ---
\end{verbatim}
Depending on the base LLM, the visual input was provided either as a video URI or as a set of key frames; the placeholder \verb+{{VIDEO\FRAMES}}+ denotes this modality.
The CoT variants added a \texttt{REASONING FRAMEWORK} section to the system prompt to enforce a step-by-step thought process. The multimodal variants included placeholders for all four modalities and a protocol requiring a review of each.

\subsection{Inference Configuration}
\label{app:inference}

The following fixed decoding parameters were used for all API calls:
\begin{itemize}
    \item Decoding: \texttt{temperature=0}, \texttt{top-p=1}, \texttt{n=1}.
    \item Max output tokens: \texttt{8096}.
    \item Log-probabilities: top-20 per token.
\end{itemize}
Tokens were segmented into classification and reasoning spans based on the structured output. Assets (images, videos) were passed as Google Cloud Storage URIs, and malformed JSON responses (fewer than 2\% of total requests) were retried up to three times.
